%\section{Discussion: Remove parts into experiments and rest in conclusion TODO}
%\gls{gem} matches or exceed baseline models for molecular graph benchmarks in both the unconditional and conditional setting. It enables compositional, inference-time control by adding energy terms, as shown in the conditional
%experiments (Section~\ref{sec:experiments}). The learned energy also enables interpretable diagnostics,
%including geodesic analysis and likelihood-ratio calibration.

%\paragraph{Limitations} At very small inference budgets, diffusion-based baselines can reach strong performance with fewer steps;
%Figure~\ref{fig:moses-vun-steps} shows that in the few-step regime our current \gls{gem} sampler trails the
%diffusion baseline. More broadly, \glspl{ebm} for generation rely on gradient-based velocity parameterization,
%which typically requires both forward and backward passes per step and implies an inherent $\sim$2--3$\times$
%theoretical slowdown relative to purely forward samplers.


%\paragraph{Future directions} 
%A promising extension is 3D molecular generation: GEM can retain a strong graph prior while adding geometric energies such as distance/torsion constraints or pharmacophore terms at inference time. Additionally, GEM’s explicit energy landscape could be used for a broader diagnostics toolkit: relative likelihood rankings and
%energy-weighted geodesics can be used to quantify generalization or flag out-of-distribution regions. In particular, if the minimum-energy geodesic between two molecular (scaffold) families has a high-energy barrier, the model likely has low support (is effectively out-of-distribution) in the intermediates, so generations requiring that edit transition should be treated as low-trust.


%Old
%An immediate extension is 3D molecule generation, composing a strong graph prior with geometric
%constraints at inference time. Relative likelihoods can support scalable refinement and
%generalization diagnostics by testing energy rankings on held-out regions.

\section{Conclusion}
We introduced \emph{Graph Energy Matching} (GEM), a framework motivated by the map-optimization perspective of the \gls{jko} scheme, extending continuous energy matching principles to discrete graph spaces. GEM integrates deterministic, transport-aligned graph edits guiding sampling toward high-probability regions, with principled \gls{mcmc} mixing under a unified, scalar energy model, bridging graph generation and discrete energy-based modeling. This design provides two practical advantages: (i) scalable, high-quality unconditional molecular generation from both noise-initialized and data-driven (editing) initializations, and (ii) flexible inference-time conditional generation together with a learned relative-likelihood structure, enabling downstream tasks such as geodesic interpolation. Empirically, GEM is on pair or better against strong diffusion/\gls{ctmc} baselines on MOSES and QM9, and it performs effective property optimization under threshold constraints. Beyond generation quality, GEM exposes an explicit energy landscape that can be utilized for relative likelihood comparisons and geodesic analyses, providing a complementary diagnostic view of model support and controllability in learned molecular space. 

\paragraph{Limitations}: Broadly, \glspl{ebm} for generation rely on gradient-based velocity parameterization,
which typically requires both forward and backward passes per step and implies an inherent $\sim$2--3$\times$
theoretical slowdown relative to purely forward samplers. While gradient-based guidance and multi-step sampling introduce additional compute relative to purely forward generators, GEM offers a flexible and analyzable foundation for controllable graph generation.

\paragraph{Future Directions}:
A promising extension is 3D molecular generation: GEM can retain a strong graph prior while adding geometric energies such as distance or torsion constraints at inference time. Additionally, GEM’s explicit energy landscape offers a valuable diagnostics toolkit: by comparing energy profiles across datasets, one can directly quantify generalization to unseen distributions or flag out-of-distribution regions.